% !TeX root = ../main.tex
\section{Lambda calculus}
Lambda calculus is a formal system that serves as the theoretical foundation for functional programming languages. The core idea is surprisingly simple: everything reduces to functions, how they are defined and how they are applied. This is the untyped calculus; once we add a type discipline on top we get the simply typed lambda calculus of the next section.

\subsection{basic syntax}
There are only three constructs in the language:


\begin{itemize}
    \item \textbf{Variables:} $x, y, z \dots$
    \item \textbf{Abstraction:} $\lambda x.\, e$ — defines a function with
    parameter $x$ and body $e$.
    \item \textbf{Application:} $e_1\, e_2$ — applies function $e_1$ to
    argument $e_2$.
\end{itemize}

from just these three pieces, any computation can be expressed.


\subsection{Free and Bound Variables}
In $\lambda x.\, M$, we say $x$ is \textbf{bound}: the abstraction defines
its scope. A variable that appears outside the scope of any $\lambda$ is
called \textbf{free}.

\textbf{$\alpha$-equivalence} tells us that the name of a bound variable
does not matter, as long as renaming causes no conflicts:
\[
  \lambda x.\, x \;\equiv_\alpha\; \lambda y.\, y
\]

\subsection{Reduction Rules}
Computation happens through two main rules:
\begin{itemize}
    \item \textbf{$\beta$-reduction} — substitutes the argument into the
    function body:
    \[(\lambda x.\, e)\, v \;\to\; e[x := v]\]
    This is the fundamental computational step.

    \item \textbf{$\eta$-reduction} — two functions are equal if they produce
    the same result for every argument:
    \[\lambda x.\, (f\, x) \;\to\; f \qquad \text{(if $x$ is not free in
    $f$)}\]
\end{itemize}

\subsection{Evaluation Strategies}
In what order do we reduce? There are several options, each with different
consequences:

\begin{enumerate}
    \item \textbf{Normal Order:} always reduces the leftmost, outermost redex
    first. Guaranteed to find a normal form if one exists.
    \item \textbf{Call-by-Value:} evaluates the argument before passing it
    to the function (used in C, Java, Python).
    \item \textbf{Call-by-Name:} passes the argument unevaluated — the basis
    for lazy evaluation.
\end{enumerate}

\subsection{Church Encodings}
Since lambda calculus has no built-in types or literals, we need to
\emph{encode} familiar values as functions:

\begin{itemize}
    \item \textbf{Booleans:}
    \begin{align*}
        \mathit{TRUE}  &\equiv \lambda t.\, \lambda f.\, t \\
        \mathit{FALSE} &\equiv \lambda t.\, \lambda f.\, f \\
        \mathit{AND}   &\equiv \lambda p.\, \lambda q.\, p\, q\, p
    \end{align*}
    \item \textbf{Church Numerals} — the number $n$ applies $f$ exactly
    $n$ times:
    \begin{align*}
        0    &\equiv \lambda f.\, \lambda x.\, x \\
        1    &\equiv \lambda f.\, \lambda x.\, f\, x \\
        2    &\equiv \lambda f.\, \lambda x.\, f\,(f\, x) \\
        \mathit{SUCC} &\equiv \lambda n.\, \lambda f.\, \lambda x.\,
                               f\,(n\, f\, x)
    \end{align*}
\end{itemize}

\subsection{Recursion: The Y Combinator}

Lambda functions are anonymous, so recursion is not straightforward: there is
no name to refer to inside the body. The solution is Curry's
\textbf{fixed-point combinator}:
\[
  Y = \lambda f.\,(\lambda x.\, f\,(x\, x))\,(\lambda x.\, f\,(x\, x))
\]
It satisfies $Y\, g = g\,(Y\, g)$, which lets us build recursive functions
without ever naming them: $Y\,g$ unfolds itself exactly as many times as $g$
needs.

\subsection{Exercises}
\begin{itemize}
  \item Reduce $(\lambda x.\, x\, x)\,(\lambda y.\, y)$ to normal form,
  showing every $\beta$-reduction step.
  \item Using the Church encoding of booleans, define $\mathit{OR}$ and
  $\mathit{NOT}$, and check that $\mathit{OR}\;\mathit{TRUE}\;\mathit{FALSE}$
  reduces to $\mathit{TRUE}$.
  \item Give the Church numeral for $3$ and show that
  $\mathit{SUCC}\;2 \equiv_\beta 3$.
  \item Verify the defining property of the Y combinator: unfold
  $Y\,g$ one step and check it $\beta$-reduces to $g\,(Y\,g)$.
\end{itemize}